
\documentclass{article}
\usepackage{colortbl}
\usepackage{makecell}
\usepackage{multirow}
\usepackage{supertabular}

\begin{document}

\newcounter{utterance}

\centering \large Interaction Transcript for game `cladder', experiment `full\_v1.5\_default', episode 525 with Qwen3.5{-}4B.
\vspace{24pt}

{ \footnotesize  \setcounter{utterance}{1}
\setlength{\tabcolsep}{0pt}
\begin{supertabular}{c@{$\;$}|p{.15\linewidth}@{}p{.15\linewidth}p{.15\linewidth}p{.15\linewidth}p{.15\linewidth}p{.15\linewidth}}
    \# & \multicolumn{2}{c}{Player} && \multicolumn{2}{c}{Game Master} \\
    \hline

    \theutterance \stepcounter{utterance}  
    & & & \multicolumn{4}{p{0.6\linewidth}}{
        \cellcolor[rgb]{0.9,0.9,0.9}{
            \makecell[{{p{\linewidth}}}]{
                \texttt{\tiny{[P1$\langle$GM]}}
                \texttt{You are an expert in causal inference. The following question is not a typical commonsense query, but rather a meticulously designed question created by a professor specializing in causal inference, intended to assess the students' mastery of the course content.} \\
\texttt{Imagine a self{-}contained, hypothetical world with only the following conditions, and without any unmentioned factors or causal relationships: Unobserved confounders has a direct effect on drug taken and cholesterol level. Treatment assignment has a direct effect on drug taken. Drug taken has a direct effect on cholesterol level. Unobserved confounders is unobserved. For patients not assigned the drug treatment, the probability of low cholesterol is 63\%. For patients assigned the drug treatment, the probability of low cholesterol is 52\%. For patients not assigned the drug treatment, the probability of taking of all assigned drugs is 37\%. For patients assigned the drug treatment, the probability of taking of all assigned drugs is 74\%. Will taking of all assigned drugs increase the chance of low cholesterol?} \\
\texttt{Start your answer with 'Yes' or 'No', followed by additional reasoning or evidence to support your explanation.} \\
            }
        }
    }
     \\ \\

    \theutterance \stepcounter{utterance}  
    & \multicolumn{4}{p{0.6\linewidth}}{
        \cellcolor[rgb]{0.561,0.737,0.561}{
            \makecell[{{p{\linewidth}}}]{
                \texttt{\tiny{[P1$\rangle$GM]}}
                \texttt{No} \\
\\ 
\texttt{**Reasoning:**} \\
\\ 
\texttt{To determine whether taking the drug increases the chance of low cholesterol, we must calculate the **Average Treatment Effect on the Treated (ATT)**, which represents the causal effect of the drug on the subgroup of patients who actually took it. The question asks about the effect of "taking of all assigned drugs," implying we are looking at the counterfactual comparison for those who took the drug: *What would have happened to them if they had not taken the drug, versus what actually happened?*} \\
\\ 
\texttt{We are given the following data points in this hypothetical world:} \\
\texttt{1.  **Outcome for Non{-}Treated (Control Group):** \$P(\textbackslash{}text\{Low Cholesterol\} | \textbackslash{}text\{Not Treated\}) = 0.63\$.} \\
\texttt{2.  **Outcome for Treated (Actual Outcome):** \$P(\textbackslash{}text\{Low Cholesterol\} | \textbackslash{}text\{Treated\}) = 0.52\$.} \\
\texttt{3.  **Treatment Compliance:**} \\
\texttt{    *   \$P(\textbackslash{}text\{Took Drug\} | \textbackslash{}text\{Not Assigned\}) = 0.37\$} \\
\texttt{    *   \$P(\textbackslash{}text\{Took Drug\} | \textbackslash{}text\{Assigned\}) = 0.74\$} \\
\\ 
\texttt{**Step 1: Calculate the Probability of Taking the Drug for the Treated Group**} \\
\texttt{The group "Assigned the drug treatment" consists of two subgroups: those who took the drug and those who did not.} \\
\texttt{\$\$P(\textbackslash{}text\{Took\} | \textbackslash{}text\{Assigned\}) = P(\textbackslash{}text\{Took\} | \textbackslash{}text\{Assigned\}) \textbackslash{}times 1 = 0.74\$\$} \\
\texttt{\$\$P(\textbackslash{}text\{Did Not Take\} | \textbackslash{}text\{Assigned\}) = 1 {-} 0.74 = 0.26\$\$} \\
\\ 
\texttt{**Step 2: Calculate the Counterfactual Outcome for the Treated Group**} \\
\texttt{To find the causal effect, we need to estimate what the outcome would have been for the 26\% of the treated group who *did not* take the drug (the compliers/non{-}compliers within the assigned group). Since the prompt states "without any unmentioned factors," we assume that for those who did not take the drug, their outcome is determined by the baseline probability of the non{-}treated group (as there are no other unobserved confounders affecting the outcome directly other than the drug and the assignment mechanism which is already accounted for in the group definition).} \\
\\ 
\texttt{*Note: In standard causal inference problems of this type, if a patient is assigned but does not take the treatment, and no other specific factors are mentioned, we treat their potential outcome under "No Treatment" as the same as the observed outcome of the non{-}treated population.*} \\
\\ 
\texttt{\$\$P(\textbackslash{}text\{Low Cholesterol\} | \textbackslash{}text\{Assigned, Did Not Take\}) = P(\textbackslash{}text\{Low Cholesterol\} | \textbackslash{}text\{Not Treated\}) = 0.} \\
            }
        }
    }
    & & \\ \\

    \theutterance \stepcounter{utterance}  
    & & & \multicolumn{2}{p{0.3\linewidth}}{
        \cellcolor[rgb]{0.9,0.9,0.9}{
            \makecell[{{p{\linewidth}}}]{
                \texttt{\tiny{[GM$|$GM]}}
                \texttt{No} \\
            }
        }
    }
    & & \\ \\

    \theutterance \stepcounter{utterance}  
    & & & \multicolumn{2}{p{0.3\linewidth}}{
        \cellcolor[rgb]{0.9,0.9,0.9}{
            \makecell[{{p{\linewidth}}}]{
                \texttt{\tiny{[GM$|$GM]}}
                \texttt{no} \\
            }
        }
    }
    & & \\ \\

    \theutterance \stepcounter{utterance}  
    & & & \multicolumn{2}{p{0.3\linewidth}}{
        \cellcolor[rgb]{0.9,0.9,0.9}{
            \makecell[{{p{\linewidth}}}]{
                \texttt{\tiny{[GM$|$GM]}}
                \texttt{game\_result = WIN} \\
            }
        }
    }
    & & \\ \\

\end{supertabular}
}

\end{document}
